% =========================================================================
\section{Cálculo integral vectorial: los teoremas integrales fundamentales}\label{sec:integral}
% =========================================================================

% -------------------------------------------------------------------------
\subsection{Del teorema fundamental del cálculo a los teoremas vectoriales}\label{sec:tfc}
% -------------------------------------------------------------------------
Recordemos el teorema fundamental del cálculo en una variable: para una
función $f$ derivable en $[a,b]$,
\begin{equation}\label{eq:tfc}
  \int_{a}^{b}\frac{\dif f}{\dif x}\,\dif x=f(b)-f(a).
\end{equation}
La idea es la de una \emph{suma de variaciones}: si dividimos $[a,b]$ en
intervalitos de longitud $\dif x$, en cada uno de ellos la función cambia en
$\dif f$; al sumar todos los cambios infinitesimales, las contribuciones
intermedias se cancelan por parejas y solo queda la diferencia entre lo que
vale la función en los extremos (\cref{fig:tfc}). Es la primera instancia de un
principio que se repite en todo el cálculo vectorial:

\begin{quote}\itshape
La integral de una «derivada» sobre una región es igual a la integral de la
función sobre el borde de esa región.
\end{quote}

\begin{figure}[ht]
\centering
\begin{tikzpicture}[>=Stealth, x=1cm, y=1cm]
  \draw[eje] (-0.3,0) -- (8.6,0) node[right, text=black] {$x$};
  \draw[eje] (0,-0.3) -- (0,4.4) node[above, text=black] {$f(x)$};
  \def\fx#1{0.9+0.55*(#1)+0.9*sin(deg(0.55*(#1)))}
  % barras de variación
  \foreach \xx in {1.0,1.6,...,7.0}{%
    \pgfmathsetmacro{\ya}{\fx{\xx}}\pgfmathsetmacro{\yb}{\fx{\xx+0.6}}%
    \draw[azulclaro, line width=0.7pt] (\xx,\ya) -- (\xx+0.6,\ya);
    \draw[rojoteo!70, line width=1.0pt] (\xx+0.6,\ya) -- (\xx+0.6,\yb);}
  \draw[azulprof, line width=1.3pt, domain=1.0:7.6, samples=120, smooth]
    plot (\x, {\fx{\x}});
  \foreach \xx/\lab in {1.0/a,7.6/b}{%
    \pgfmathsetmacro{\yy}{\fx{\xx}}%
    \draw[densely dashed, grispizarra] (\xx,0) -- (\xx,\yy);
    \fill[azulprof] (\xx,\yy) circle (2pt);
    \node[below] at (\xx,0) {$\lab$};}
  \node[font=\small, text=rojoteo, anchor=west] at (3.0,3.7) {$\dif f$};
  \draw[->, rojoteo, line width=0.5pt] (3.0,3.65) -- (2.55,3.0);
  \node[font=\small, text=azulacero, anchor=north] at (3.9,0.35) {$\dif x$};
  \draw[|<->|, azulacero, line width=0.5pt] (3.4,0.3) -- (4.0,0.3);
  \node[font=\small, text=azulprof, rotate=90, anchor=south] at (9.1,{(\fx{1.0}+\fx{7.6})/2}) {$f(b)-f(a)=\textstyle\sum\dif f$};
  \draw[|<->|, grispizarra, line width=0.5pt] (8.5,{\fx{1.0}}) -- (8.5,{\fx{7.6}});
  \draw[densely dashed, grispizarra] (7.6,{\fx{7.6}}) -- (8.5,{\fx{7.6}});
  \draw[densely dashed, grispizarra] (1.0,{\fx{1.0}}) -- (8.5,{\fx{1.0}});
\end{tikzpicture}
\caption{Teorema fundamental del cálculo como suma de variaciones: la suma de
los incrementos $\dif f$ (en rojo) sobre una partición de $[a,b]$ es el
incremento total $f(b)-f(a)$.}
\label{fig:tfc}
\end{figure}

El «borde» de un intervalo $[a,b]$ son sus dos extremos, a los que se asigna una
orientación: $b$ cuenta con signo $+$ (punto de salida) y $a$ con signo $-$
(punto de entrada). Al aumentar la dimensión del dominio de integración se
obtienen tres generalizaciones, que son los teoremas de este apartado:
\begin{center}
\renewcommand{\arraystretch}{1.3}
\begin{tabular}{@{}llll@{}}
\toprule
Dominio $\Omega$ & Dimensión & «Derivada» & Borde $\partial\Omega$\\
\midrule
curva $C$ & 1 & gradiente $\nab f$ & dos puntos $\{\vc{a}^{-},\vc{b}^{+}\}$\\
superficie $S$ & 2 & rotacional $\nab\times\vc{A}$ & curva cerrada $L=\partial S$\\
sólido $V$ & 3 & divergencia $\nab\cdot\vc{A}$ & superficie cerrada $S=\partial V$\\
\bottomrule
\end{tabular}
\end{center}

% -------------------------------------------------------------------------
\subsection{Integral de línea}\label{sec:linea}
% -------------------------------------------------------------------------
\begin{definicion}[Integral de línea]\label{def:linea}
Sean
\begin{itemize}
  \item $C\subseteq\R^{n}$ una curva regular parametrizada por una aplicación
        $\gamma:[a,b]\to\R^{n}$ de clase $C^{1}$ con $\gamma'(t)\neq\vc{0}$, y
  \item $\vc{F}:C\to\R^{n}$ un campo vectorial continuo sobre la curva.
\end{itemize}
Se define la \emph{integral de línea} (o \emph{de camino}, o \emph{de trabajo})
de $\vc{F}$ a lo largo de $C$ como la integral ordinaria
\begin{equation}\label{eq:linea}
  \int_{C}\vc{F}\cdot\dl\defeq\int_{a}^{b}\bigl\langle\vc{F}(\gamma(t)),\gamma'(t)\bigr\rangle\,\dif t
  =\int_{a}^{b}\bigl[(\vc{F}\circ\gamma)\cdot\gamma'\bigr](t)\,\dif t .
\end{equation}
El símbolo $\dl\defeq\gamma'(t)\,\dif t$ es el \emph{elemento de longitud
dirigido}: un desplazamiento infinitesimal tangente a la curva,
$\dl=\sum_i\dif x_i\,\uh{i}$ con $\dif x_i=x_i'(t)\,\dif t$. Si $C$ es cerrada se
escribe $\oint_C$ y se habla de \emph{circulación}.
\end{definicion}

Físicamente, $\int_C\vc{F}\cdot\dl$ es el \emph{trabajo} que realiza el campo de
fuerzas $\vc{F}$ sobre una partícula que recorre $C$: en cada tramo infinitesimal
solo cuenta la componente de $\vc{F}$ a lo largo del desplazamiento.

\begin{notacion}[Procedimiento para calcular una integral de línea]\label{not:procedimiento-linea}
Se procede en cinco pasos.\begin{enumerate}[label=\textbf{\arabic*.}, itemsep=1pt]
  \item \emph{Parametrizar} la curva: $\vc{r}=\gamma(t)=(x_1(t),\ldots,x_n(t))$, $t\in[a,b]$.
  \item \emph{Derivar}: calcular el vector tangente $\gamma'(t)=(x_1'(t),\ldots,x_n'(t))$.
  \item \emph{Restringir} el campo a la trayectoria: $(\vc{F}\circ\gamma)(t)=\vc{F}(\gamma(t))$.
  \item \emph{Proyectar}: producto escalar $\vc{F}(\gamma(t))\cdot\gamma'(t)$, que es una función de una sola variable.
  \item \emph{Integrar} esa función de $t$ sobre $[a,b]$.
\end{enumerate}
\end{notacion}

\paragraph{En coordenadas.} Si $\vc{F}=(f_1,\ldots,f_n)$, la integral de línea es
la suma de $n$ integrales ordinarias:
\begin{equation}\label{eq:linea-coord}
  \int_{C}\vc{F}\cdot\dl=\int_{a}^{b}\sum_{k=1}^{n}(f_k\circ\gamma)(t)\,x_k'(t)\,\dif t
  =\sum_{k=1}^{n}\int_{a}^{b}f_k\bigl(\gamma(t)\bigr)\,x_k'(t)\,\dif t
\end{equation}
(aditividad de la integral).

\begin{definicion}[Notación diferencial y 1-formas]\label{def:diferencial-linea}
Con el cambio de variable $\dif x_k=x_k'(t)\,\dif t$ se tiene
$f_k(\vc{r})\,\dif x_k=f_k(x_1(t),\ldots,x_n(t))\,x_k'(t)\,\dif t$ y la integral de
línea queda
\[
  \int_{C}\vc{F}\cdot\dl=\int_{C}f_1\,\dif x_1+f_2\,\dif x_2+\cdots+f_n\,\dif x_n ,
\]
donde cada sumando se integra por separado. A una expresión de este tipo,
$\omega=f_1\,\dif x_1+\cdots+f_n\,\dif x_n=\sum_i f_i\,\dif x_i$ con $f_k$
campos escalares, se le llama \emph{1-forma diferencial}.
\end{definicion}

En dos y tres dimensiones, con $\vc{F}=(P,Q)$ y $\vc{F}=(P,Q,R)$ respectivamente,
\begin{equation}\label{eq:linea-2d3d}
  \int_{C}\vc{F}\cdot\dl=\int_{C}P\,\dif x+Q\,\dif y
  \qquad\text{y}\qquad
  \int_{C}\vc{F}\cdot\dl=\int_{C}P\,\dif x+Q\,\dif y+R\,\dif z,
\end{equation}
donde, por ejemplo, $\int_{C}P\,\dif x=\int_a^bP\bigl(x(t),y(t),z(t)\bigr)\,x'(t)\,\dif t$.

\begin{ejemplo}[Segmento rectilíneo]\label{ej:segmento}
Sea $C=[\vc{p},\vc{q}]$ el segmento que une dos puntos de $\R^{n}$. Su
parametrización estándar es $\sigma(t)=\vc{p}+t\,\vc{h}$ con
$\vc{h}=\vc{q}-\vc{p}$ y $t\in[0,1]$: cumple $\sigma(0)=\vc{p}$,
$\sigma(1)=\vc{q}$ y $\sigma'(t)=\vc{h}$ constante. Para cualquier campo
continuo,
\[
  \int_{[\vc{p},\vc{q}]}\vc{F}\cdot\dl
  =\int_{0}^{1}\vc{F}(\vc{p}+t\,\vc{h})\cdot\vc{h}\,\dif t
  =\int_{0}^{1}\vc{F}\bigl(\vc{p}+t(\vc{q}-\vc{p})\bigr)\cdot(\vc{q}-\vc{p})\,\dif t .
\]
\end{ejemplo}

\begin{proposicion}[Propiedades de la integral de línea]\label{prop:prop-linea}
\begin{enumerate}[label=(\alph*)]
  \item \emph{Linealidad:} $\displaystyle\int_C(\alpha\vc{F}+\beta\vc{G})\cdot\dl=\alpha\int_C\vc{F}\cdot\dl+\beta\int_C\vc{G}\cdot\dl$.
  \item \emph{Aditividad:} si $C=C_1\uplus C_2$ (la segunda curva empieza donde acaba la primera), $\int_C=\int_{C_1}+\int_{C_2}$.
  \item \emph{Cambio de orientación:} si $C^{\mathrm{op}}$ es la misma curva recorrida en sentido contrario, $\int_{C^{\mathrm{op}}}\vc{F}\cdot\dl=-\int_C\vc{F}\cdot\dl$.
  \item \emph{Invariancia por reparametrización:} si $\alpha:[a,b]\to C$ y
        $\beta:[A,B]\to C$ son dos parametrizaciones regulares de la misma curva y
        recorren $C$ en el mismo sentido, las integrales \eqref{eq:linea}
        coinciden.
\end{enumerate}
\end{proposicion}
\begin{proof}
Las tres primeras se siguen de las propiedades de la integral de Riemann.
Para (d), sea $\varphi=\beta^{-1}\circ\alpha:[a,b]\to[A,B]$, un cambio de
variable de clase $C^{1}$ y biyectivo, de modo que $\alpha=\beta\circ\varphi$ y
$\alpha'(t)=\beta'(\varphi(t))\,\varphi'(t)$. Entonces
\[
  \int_a^b\vc{F}(\alpha(t))\cdot\alpha'(t)\,\dif t
  =\int_a^b\vc{F}\bigl(\beta(\varphi(t))\bigr)\cdot\beta'(\varphi(t))\,\varphi'(t)\,\dif t
  =\int_A^B\vc{F}(\beta(s))\cdot\beta'(s)\,\dif s ,
\]
con $s=\varphi(t)$, siempre que $\varphi'>0$ (mismo sentido). Si $\varphi'<0$, los
límites de integración se invierten y la integral cambia de signo, que es (c).
\end{proof}

La integral de línea está bien definida: depende de la curva y de su
orientación, pero no de cómo se parametrice.

% -------------------------------------------------------------------------
\subsection{El teorema del gradiente}\label{sec:teorema-gradiente}
% -------------------------------------------------------------------------
\begin{teorema}[Teorema del gradiente o fundamental de las integrales de línea]\label{teo:gradiente}
Sea $f$ un campo escalar de clase $C^{1}$ en una región $W\subseteq\R^{n}$ y sea
$C\subset W$ un arco regular orientado cuyo borde es $\partial C=\{\vc{a}^{-},\vc{b}^{+}\}$
(punto de partida $\vc{a}$, punto de llegada $\vc{b}$). Entonces
\begin{equation}\label{eq:teo-gradiente}
  \boxed{\;\int_{C}\nab f\cdot\dl=f(\vc{b})-f(\vc{a})\;}
\end{equation}
\end{teorema}
\begin{proof}
Sea $\gamma:[a,b]\to C$ una parametrización compatible con la orientación, con
$\gamma(a)=\vc{a}$ y $\gamma(b)=\vc{b}$. Por la \cref{prop:cadena},
$\nab f(\gamma(t))\cdot\gamma'(t)=(f\circ\gamma)'(t)$, y por el teorema
fundamental del cálculo en una variable, \eqref{eq:tfc},
\[
  \int_C\nab f\cdot\dl=\int_a^b\nab f(\gamma(t))\cdot\gamma'(t)\,\dif t
  =\int_a^b(f\circ\gamma)'(t)\,\dif t=f(\gamma(b))-f(\gamma(a))=f(\vc{b})-f(\vc{a}).\qedhere
\]
\end{proof}

El valor de la integral queda determinado única y exclusivamente por la
diferencia de potencial en los extremos del camino. Con la notación de la
\cref{sec:tfc}: la integral de la «derivada» $\nab f$ sobre $C$ es igual a $f$
evaluada en el borde $\partial C$, contando $+f(\vc{b})$ y $-f(\vc{a})$. Dos
consecuencias inmediatas son de gran importancia en física.

\begin{corolario}[Independencia del camino y circulación nula]\label{cor:gradiente}
Sea $\vc{F}=\nab f$ un campo gradiente en $W$.
\begin{enumerate}[label=(\alph*)]
  \item El trabajo de $\vc{F}$ entre dos puntos es el mismo sobre cualquier curva
        que los una: depende solo de los extremos, no del camino.
  \item Para toda curva cerrada $C\subset W$ (con $\vc{r}(a)=\vc{r}(b)$),
        \begin{equation}\label{eq:circulacion-nula}
          \oint_C\nab f\cdot\dl=0 .
        \end{equation}
\end{enumerate}
\end{corolario}
\begin{proof}
(a) Es inmediato de la \cref{eq:teo-gradiente}. (b) Si la curva es cerrada,
$\vc{b}=\vc{a}$ y $f(\vc{b})-f(\vc{a})=0$.
\end{proof}

\begin{definicion}[Independencia del camino y circulación nula]\label{def:indep-camino}
Un campo $\vc{F}$ en $W$ tiene la propiedad de \emph{independencia del camino}
si $\int_\alpha\vc{F}\cdot\dl=\int_\beta\vc{F}\cdot\dl$ para cualquier par de
curvas $\alpha,\beta\subset W$ con los mismos extremos orientados; y tiene
\emph{circulación nula} si $\oint_C\vc{F}\cdot\dl=0$ para toda curva cerrada
$C\subset W$.
\end{definicion}

\begin{proposicion}\label{prop:equiv-indep}
Un campo continuo tiene independencia del camino en $W$ si y solo si tiene
circulación nula en $W$.
\end{proposicion}
\begin{proof}
($\Leftarrow$) Si $A$ y $B$ unen los mismos extremos, $C=A\uplus B^{\mathrm{op}}$ es
un circuito, y por la \cref{prop:prop-linea},
$0=\oint_C\vc{F}\cdot\dl=\int_A\vc{F}\cdot\dl-\int_B\vc{F}\cdot\dl$.
($\Rightarrow$) Si $C$ es cerrada, elegimos dos puntos distintos de $C$, que la
dividen en dos arcos $A$ y $B$ con los mismos extremos y $C=A\uplus B^{\mathrm{op}}$;
entonces $\oint_C\vc{F}\cdot\dl=\int_A\vc{F}\cdot\dl-\int_B\vc{F}\cdot\dl=0$.
\end{proof}

\begin{proposicion}[Recíproco del teorema del gradiente]\label{prop:reciproco}
Un campo continuo $\vc{F}$ en una región $W$ con independencia del camino es un
campo gradiente: fijado $\vc{p}\in W$, la función
$f(\vc{r})=\int_{\vc{p}}^{\vc{r}}\vc{F}\cdot\dl$ (calculada sobre cualquier
curva de $W$ que una $\vc{p}$ con $\vc{r}$) cumple $\nab f=\vc{F}$.
\end{proposicion}
\begin{proof}
Por independencia del camino, $f$ está bien definida. Para $h$ pequeño,
$f(\vc{r}+h\,\uh{i})-f(\vc{r})$ es la integral de $\vc{F}$ sobre el segmento que
une $\vc{r}$ con $\vc{r}+h\,\uh{i}$, es decir,
$\int_0^hF_i(\vc{r}+t\,\uh{i})\,\dif t$. Dividiendo por $h$ y haciendo
$h\to0$, la continuidad de $F_i$ da $\partial_if(\vc{r})=F_i(\vc{r})$.
\end{proof}

Reuniendo todo: para campos continuos, \emph{ser gradiente $\Leftrightarrow$
independencia del camino $\Leftrightarrow$ circulación nula}; y ser irrotacional es
condición necesaria (\cref{cor:grad-irrot}), y suficiente en regiones
estrelladas.

% Figura: teorema del gradiente. (a) Trabajo de un gradiente entre a y b por
% dos caminos; (b) circulación nula en un circuito cerrado.
\newcommand{\curvasNivelElipt}{%
  \foreach \c/\pc in {0.6/100,1.2/80,1.8/60,2.4/40}{%
    \draw[azulprof!\pc!azulclaro, line width=0.8pt] (0,0) ellipse ({2*\c} and {\c});}
  \fill[azulprof] (0,0) circle (1.4pt);}
\begin{figure}[ht]
\centering
\begin{minipage}[t]{0.49\textwidth}\centering
\begin{tikzpicture}[>=Stealth, x=0.74cm, y=0.74cm, baseline=(current bounding box.north)]
  \clip (-5.1,-3.2) rectangle (5.1,3.2);
  \curvasNivelElipt
  % camino principal: x(t)=-4.2+6t, y(t)=-2.2+3.0t-1.4 sin(pi t)
  \draw[trayecto, domain=0:1, samples=60, smooth, variable=\t]
    plot ({-4.0+5.8*\t},{-2.0+3.0*\t-1.2*sin(180*\t)});
  % camino alternativo: segmento
  \draw[acento!60, densely dashed, line width=0.8pt] (-4.0,-2.0) -- (1.8,1.0);
  % vectores en puntos del camino principal
  \foreach \t in {0.18,0.52,0.86}{%
    \pgfmathsetmacro{\px}{-4.0+5.8*\t}\pgfmathsetmacro{\py}{-2.0+3.0*\t-1.2*sin(180*\t)}%
    \pgfmathsetmacro{\dx}{5.8}\pgfmathsetmacro{\dy}{3.0-1.2*pi*cos(180*\t)}%
    \pgfmathsetmacro{\dn}{sqrt(\dx*\dx+\dy*\dy)}%
    \pgfmathsetmacro{\gx}{-\px/2}\pgfmathsetmacro{\gy}{-2*\py}%
    \pgfmathsetmacro{\gn}{sqrt(\gx*\gx+\gy*\gy)}%
    \draw[->, azulprof, line width=1pt] (\px,\py) -- ++({0.8*\dx/\dn},{0.8*\dy/\dn});
    \draw[->, line width=0.9pt, cm2] (\px,\py) -- ++({1.2*\gx/\gn},{1.2*\gy/\gn});
    \fill (\px,\py) circle (0.8pt);}
  \fill[acento] (-4.0,-2.0) circle (1.2pt) node[below left=0pt, text=black] {$\vc{a}$};
  \fill[acento] (1.8,1.0) circle (1.2pt) node[above right=0pt, text=black] {$\vc{b}$};
  \node[font=\small, text=acento] at (-0.9,-2.1) {$C$};
  
  \node[fill=white, fill opacity=0.85, text opacity=1, inner sep=2pt, font=\footnotesize, anchor=south west]
       at (-5.0,2.45) {\textcolor{azulprof}{$\to\ \dl$}\quad\textcolor{cm2}{$\to\ \nab f$}};
\end{tikzpicture}\\[3pt]{\small\scshape (a)}\; $\displaystyle\int_C\nab f\cdot\dl=f(\vc{b})-f(\vc{a})$
\end{minipage}\hfill
\begin{minipage}[t]{0.49\textwidth}\centering
\begin{tikzpicture}[>=Stealth, x=0.74cm, y=0.74cm, baseline=(current bounding box.north)]
  \clip (-5.1,-3.2) rectangle (5.1,3.2);
  \curvasNivelElipt
  % circuito cerrado: elipse desplazada, sentido antihorario
  \def\cx{0.5}\def\cy{0.1}
  \draw[trayecto] ({\cx},{\cy}) ellipse (3.5 and 2.2);
  \foreach \s in {20,140,260}{%
    \pgfmathsetmacro{\px}{\cx+3.5*cos(\s)}\pgfmathsetmacro{\py}{\cy+2.2*sin(\s)}%
    \pgfmathsetmacro{\dx}{-3.5*sin(\s)}\pgfmathsetmacro{\dy}{2.2*cos(\s)}%
    \pgfmathsetmacro{\dn}{sqrt(\dx*\dx+\dy*\dy)}%
    \pgfmathsetmacro{\gx}{-\px/2}\pgfmathsetmacro{\gy}{-2*\py}%
    \pgfmathsetmacro{\gn}{sqrt(\gx*\gx+\gy*\gy)}%
    \draw[->, azulprof, line width=1pt] (\px,\py) -- ++({0.8*\dx/\dn},{0.8*\dy/\dn});
    \draw[->, line width=0.9pt, cm2] (\px,\py) -- ++({1.2*\gx/\gn},{1.2*\gy/\gn});
    \fill (\px,\py) circle (0.8pt);}
  \node[font=\small, text=acento] at (4.4,2.55) {$C$};
  \node[fill=white, fill opacity=0.85, text opacity=1, inner sep=2pt, font=\footnotesize, anchor=south west]
       at (-5.0,2.45) {\textcolor{azulprof}{$\to\ \dl$}\quad\textcolor{cm2}{$\to\ \nab f$}};
\end{tikzpicture}\\[3pt]{\small\scshape (b)}\; $\displaystyle\oint_C\nab f\cdot\dl=0$
\end{minipage}
\caption{Teorema del gradiente sobre el mapa de curvas de nivel de una
colina $f$ (ahora la cima es el máximo). (a)~Sobre la curva $C$ (y sobre
cualquier otra con los mismos extremos, p.\,ej. el segmento discontinuo) la
integral suma, en cada tramo, la proyección del gradiente sobre el desplazamiento
$\dl$: tramos cuesta arriba suman, tramos cuesta abajo restan, y el balance es
solo $f(\vc{b})-f(\vc{a})$. (b)~En un circuito cerrado, lo que se sube se baja:
la circulación es cero.}
\label{fig:teor_grad}
\end{figure}


% -------------------------------------------------------------------------
\subsection{Integral de flujo}\label{sec:flujo}
% -------------------------------------------------------------------------
Pasamos a la integral que formaliza la noción de \emph{flujo} de un campo a
través de una superficie en $\R^{3}$. Ahora el vector relevante es la normal.

\begin{definicion}[Superficie parametrizada y vector normal]\label{def:superficie}
Una \emph{superficie parametrizada} $S\subset\R^{3}$ es la imagen de una
aplicación $\sigma:S^{*}\subseteq\R^{2}\to\R^{3}$, $\sigma(u,v)=(x(u,v),y(u,v),z(u,v))$,
de clase $C^{1}$ y regular ($\sigma_u\times\sigma_v\neq\vc{0}$). Su \emph{vector
normal asociado}, su normal unitaria y sus elementos de área son
\[
  \vc{n}_\sigma=\frac{\partial\sigma}{\partial u}\times\frac{\partial\sigma}{\partial v},
  \qquad
  \un=\frac{\vc{n}_\sigma}{\norm{\vc{n}_\sigma}},
  \qquad
  \dsc=\norm{\vc{n}_\sigma}\,\dif u\,\dif v,
  \qquad
  \dS=\un\,\dsc=\vc{n}_\sigma\,\dif u\,\dif v .
\]
La elección de $\un$ (de entre sus dos posibles sentidos) es una
\emph{orientación} de $S$.
\end{definicion}

\begin{definicion}[Integral de flujo]\label{def:flujo}
Dado un campo vectorial $\vc{F}$ continuo sobre una superficie orientada
$S=\sigma(S^{*})$, el \emph{flujo} de $\vc{F}$ a través de $S$ es
\begin{equation}\label{eq:flujo}
  \iint_{S}\vc{F}\cdot\dS\defeq\iint_{S^{*}}\vc{F}\bigl(\sigma(u,v)\bigr)\cdot\vc{n}_\sigma(u,v)\,\dif u\,\dif v
  =\iint_{S}\vc{F}\cdot\un\,\dsc ,
\end{equation}
donde la segunda igualdad (con $\un\,\dsc=\dS$) dice que el flujo es la integral
de \emph{área} de la componente \emph{normal} de $\vc{F}$. Para una superficie
cerrada se escribe $\oiint_S\vc{F}\cdot\dS$ y, por convenio, se toma la normal
exterior.
\end{definicion}

Solo la componente normal contribuye: un campo que corre \emph{a lo largo} de la
superficie, sin atravesarla, no aporta flujo (\cref{fig:flujo}). Ahora bien,
\emph{cualquier superficie es localmente plana}: en un entorno infinitesimal de
cada punto, el área orientada $\dS$ es un pequeño parche plano de normal $\un$,
y el flujo infinitesimal es $\vc{F}\cdot\dS=\norm{\vc{F}}\,\dsc\cos\theta$.

\begin{notacion}[Procedimiento para calcular un flujo]\label{not:procedimiento-flujo}
Se procede en seis pasos.\begin{enumerate}[label=\textbf{\arabic*.}, itemsep=1pt]
  \item \emph{Parametrizar} la superficie: $\sigma(u,v)=(x(u,v),y(u,v),z(u,v))$.
  \item \emph{Derivar}: $\sigma_u=\partial\sigma/\partial u$ y $\sigma_v=\partial\sigma/\partial v$.
  \item \emph{Normal}: $\vc{n}_\sigma=\sigma_u\times\sigma_v$ (comprobar la orientación).
  \item \emph{Restringir} el campo: $\vc{F}\circ\sigma$.
  \item \emph{Proyectar}: $(\vc{F}\circ\sigma)\cdot\vc{n}_\sigma$, en función de $u,v$.
  \item \emph{Integrar} sobre la región plana $(u,v)\in S^{*}$.
\end{enumerate}
\end{notacion}

\begin{observacion}[Notación diferencial: $2$-formas]\label{obs:2formas}
En coordenadas, $\vc{n}_\sigma=\bigl(\tfrac{\partial(y,z)}{\partial(u,v)},
\tfrac{\partial(z,x)}{\partial(u,v)},\tfrac{\partial(x,y)}{\partial(u,v)}\bigr)$,
donde $\tfrac{\partial(x,y)}{\partial(u,v)}=x_uy_v-x_vy_u$ son subdeterminantes
jacobianos. Escribiendo $\dif y\wedge\dif z\defeq\tfrac{\partial(y,z)}{\partial(u,v)}\,\dif u\,\dif v$
(y cíclicamente, $x\to y\to z\to x$), con $\vc{F}=(P,Q,R)$ queda
$\iint_S\vc{F}\cdot\dS=\iint_SP\,\dif y\wedge\dif z+Q\,\dif z\wedge\dif x+R\,\dif x\wedge\dif y$,
una \emph{2-forma diferencial}; de momento $\wedge$ es solo una notación para
estos determinantes~\cite{NAVAS}.
\end{observacion}

\begin{ejemplo}[Flujo del campo inverso del cuadrado]\label{ej:flujo-r3}
Sea $\vc{F}=\rho\,\vc{r}/r^{3}$ ($\rho$ constante, $r=\norm{\vc{r}}$) y
$S_R:x^{2}+y^{2}+z^{2}=R^{2}$. La normal unitaria a una esfera es el radio
normalizado, $\un=\vc{r}/r$, de modo que $\vc{F}\cdot\un=\rho\,r^{2}/r^{4}=\rho/R^{2}$
es constante sobre $S_R$ y
\[
  \oiint_{S_R}\vc{F}\cdot\un\,\dsc=\frac{\rho}{R^{2}}\cdot4\pi R^{2}=4\pi\rho ,
\]
independiente de $R$. Para un campo eléctrico ($\rho\propto q$) es la ley de
Gauss; su explicación profunda (la divergencia de este campo es una delta de
Dirac en el origen) se verá en la \cref{sec:delta}.
\end{ejemplo}

% Figura: flujo de un campo a través de una superficie. (a) superficie
% ondulada inmersa en un campo; (b) elemento de área dirigido y componente normal.
\begin{figure}[ht]
\centering
\begin{minipage}[t]{0.56\textwidth}\centering
\begin{tikzpicture}[baseline=(current bounding box.north)]
  \begin{axis}[hide axis, width=1.12\textwidth, height=0.95\textwidth, view={-28}{40},
      xmin=-3.2, xmax=3.2, ymin=-2.4, ymax=2.4, zmin=-1.6, zmax=2.2,
      colormap name=verdeamarillo, point meta min=-0.9, point meta max=0.9,
      /pgf/number format/use comma, clip=false]
    % campo por debajo de la superficie
    \addplot3[quiver={u=0.34, v=0.18, w=0.80}, -{Stealth[length=3pt]}, line width=0.7pt,
              acento, domain=-3.0:3.0, domain y=-2.2:2.2, samples=6, samples y=5]
      {-1.2};
    % superficie
    \addplot3[surf, shader=interp, samples=28, samples y=22, domain=-2.5:2.5, domain y=-1.8:1.8,
              opacity=0.93, mesh/ordering=x varies]
      {0.85*sin(deg(1.25*x))*cos(deg(0.5*y))+0.12*y};
    % campo por encima de la superficie
    \addplot3[quiver={u=0.34, v=0.18, w=0.80}, -{Stealth[length=3pt]}, line width=0.7pt,
              acento, domain=-3.0:3.0, domain y=-2.2:2.2, samples=6, samples y=5]
      {1.0};
    \node[font=\small, text=acento] at (axis cs:2.6,0.3,1.9) {$\vc{A}$};
    \node[font=\small, text=black!80, fill=white, fill opacity=0.7, text opacity=1, inner sep=1pt] at (axis cs:-2.0,-1.2,-0.5) {$S$};
  \end{axis}
\end{tikzpicture}\\[2pt]{\small\scshape (a)}
\end{minipage}\hfill
\begin{minipage}[t]{0.40\textwidth}\centering
\begin{tikzpicture}[>=Stealth, x=1cm, y=1cm, baseline=(current bounding box.north)]
  \shade[top color=azulclaro!55, bottom color=azulfondo, shading angle=0]
    (0,0) -- (3.4,0) -- (2.7,1.1) -- (-0.7,1.1) -- cycle;
  \draw[grispizarra!70, line width=0.4pt] (0,0) -- (3.4,0) -- (2.7,1.1) -- (-0.7,1.1) -- cycle;
  \coordinate (O) at (1.35,0.55);
  \draw[->, line width=1.6pt, azulacero] (O) -- ++(0,2.1) node[left, text=azulacero] {$\dS$};
  \draw[->, line width=1.6pt, rojoteo] (O) -- ++(2.0,1.55) node[right, text=rojoteo] {$\vc{A}$};
  % componente normal
  \draw[densely dashed, grispizarra] (3.35,2.1) -- (1.35,2.1);
  \draw[line width=3pt, azulclaro, opacity=0.9] (O) -- ++(0,1.55);
  \node[font=\footnotesize, text=azulprof, anchor=east] at (1.25,1.5) {$\vc{A}\cdot\un$};
  \draw[grispizarra, line width=0.5pt] (1.35,1.15) arc (90:37.8:0.6);
  \node[font=\small] at (1.85,1.25) {$\theta$};
  \fill (O) circle (1.4pt);
\end{tikzpicture}\\[2pt]{\small\scshape (b)}
\end{minipage}
\caption{Flujo de un campo vectorial $\vc{A}$ a través de una superficie $S$.
(a)~Un campo (flechas rojas) atraviesa una superficie curva; el flujo
contabiliza solo la parte del campo que corta a la superficie. (b)~En un
entorno de cada punto la superficie es plana: $\dS=\un\,\dsc$ es un vector de
módulo igual al área del parche y perpendicular a él, y el flujo
elemental es $\vc{A}\cdot\dS=(\vc{A}\cdot\un)\,\dsc=\norm{\vc{A}}\cos\theta\,\dsc$.}
\label{fig:flujo}
\end{figure}


% -------------------------------------------------------------------------
\subsection{El teorema de Stokes (del rotacional)}\label{sec:stokes}
% -------------------------------------------------------------------------
Una superficie $S$ con borde y su borde $\partial S$ son objetos orientados, y
sus orientaciones no son independientes.

\begin{definicion}[Orientación compatible del borde: regla de la mano derecha]\label{def:orientacion-compatible}
Dada una superficie orientada por el vector normal unitario continuo $\un$, la
\emph{orientación compatible} de su borde $\partial S$ es la del sentido de
circulación que, al acercarnos al borde, induce $\un$ según la \emph{regla del
pulgar}: con la mano derecha, el pulgar apunta según $\un$ y los demás dedos se
curvan en el sentido en que se recorre $\partial S$.
\end{definicion}

En la práctica: \emph{eligiendo el sentido en que se recorre la curva frontera
(cómo se orienta su parametrización) queda determinado el sentido de la normal}, y
viceversa.

\begin{teorema}[Teorema de Stokes]\label{teo:stokes}
Sea $\vc{A}$ un campo vectorial de clase $C^{1}$ en un entorno de una superficie
acotada y regular $S\subset\R^{3}$ cuyo borde $L=\partial S$ es una o varias
curvas cerradas. Si las orientaciones de $S$ y de $\partial S$ son compatibles,
\begin{equation}\label{eq:stokes}
  \boxed{\;\iint_{S}(\nab\times\vc{A})\cdot\dS=\oint_{\partial S}\vc{A}\cdot\dl\;}
\end{equation}
Es decir, el flujo del rotacional a través de $S$ es igual a la circulación del
campo a lo largo de su borde.
\end{teorema}
\begin{proof}[Idea de la demostración]
Sea primero $\Delta S=\Delta x\,\Delta y$ un pequeño rectángulo del plano $xy$ con
normal $\uz$ y borde recorrido en sentido antihorario. La circulación de
$\vc{A}$ alrededor de él es, a primer orden,
\begin{align*}
  \oint\vc{A}\cdot\dl&\approx A_x(x,y)\,\Delta x+A_y(x+\Delta x,y)\,\Delta y
     -A_x(x,y+\Delta y)\,\Delta x-A_y(x,y)\,\Delta y\\
  &\approx\Bigl(\pd{A_y}{x}-\pd{A_x}{y}\Bigr)\Delta x\,\Delta y
   =(\nab\times\vc{A})\cdot\uz\,\Delta S .
\end{align*}
Es decir, \emph{la circulación por unidad de área es la componente normal del
rotacional}. Una superficie cualquiera se subdivide en muchos de estos parches
infinitesimales orientados (\cref{fig:stokes}); como cada lado interior aparece
dos veces, con orientaciones opuestas, las circulaciones sobre las paredes
interiores se cancelan por parejas y solo sobrevive la circulación sobre el borde
exterior $\partial S$. Para parches curvos se cambia de variable mediante la
parametrización $\sigma$, que transforma $\uz$ en $\un$ (este es el teorema de
Green en el plano, llevado a la superficie)~\cite{NAVAS}.
\end{proof}

Para una región plana en el plano $xy$ (con $\un=\uz$), \cref{eq:stokes} se
reduce al \emph{teorema de Green}:
$\oint_{\partial S}P\,\dif x+Q\,\dif y=\iint_S\bigl(\partial_xQ-\partial_yP\bigr)\dif x\,\dif y$.

\begin{corolario}[Superficies cerradas]\label{cor:stokes-cerrada}
Si $S$ es una superficie cerrada ($\partial S=\varnothing$), el flujo del
rotacional de cualquier campo $\vc{A}$ a través de ella es nulo:
\begin{equation}\label{eq:flujo-rot-nulo}
  \oiint_{S}(\nab\times\vc{A})\cdot\dS=0 .
\end{equation}
\end{corolario}
\begin{proof}
Con $\partial S=\varnothing$ no hay circulación sobre el borde, y
\cref{eq:stokes} da $0$. El resultado es coherente con el teorema de Gauss
(\cref{teo:gauss}): si $S=\partial V$, entonces
$\oiint_S(\nab\times\vc{A})\cdot\dS=\iiint_V\nab\cdot(\nab\times\vc{A})\,\dV=0$ por
la \cref{prop:fundamentales}.
\end{proof}

Obsérvese qué ocurre en la frontera cuando $\partial S\neq\varnothing$: allí los
vórtices infinitesimales de los parches no tienen «vecino» con el que
cancelarse, por eso \emph{sí} queda la circulación sobre $L=\partial S$.

\begin{proposicion}[Rotacional y vorticidad]\label{prop:vorticidad}
La componente del rotacional según un eje unitario $\un$ en un punto $\vc{p}$ es
la circulación por unidad de área (\emph{vorticidad}) alrededor de ese eje:
\begin{equation}\label{eq:vorticidad}
  (\nab\times\vc{F})\cdot\un\,(\vc{p})=\lim_{G\to\{\vc{p}\}}\frac{1}{a(G)}\oint_{\partial G}\vc{F}\cdot\dl ,
\end{equation}
donde $G$ es un disco (o superficie similar) de área $a(G)$ contenido en un plano
ortogonal a $\un$ y que se contrae al punto $\vc{p}$.
\end{proposicion}
\begin{proof}
Por la \cref{eq:stokes}, $\frac{1}{a(G)}\oint_{\partial G}\vc{F}\cdot\dl=\frac{1}{a(G)}\iint_G(\nab\times\vc{F})\cdot\un\,\dsc$,
el valor medio de la componente normal del rotacional sobre $G$, que tiende a su
valor en $\vc{p}$ cuando $G$ se encoge a un punto (por la continuidad del
integrando).
\end{proof}

La \cref{eq:vorticidad} proporciona una definición del rotacional independiente
del sistema de coordenadas: tomando $\un=\ux,\uy,\uz$ se recuperan las tres
componentes. Con la ruedecita de paletas: la vorticidad alrededor de su eje es
proporcional a su velocidad angular de giro.

% Figura: teorema de Stokes. La circulación en los bordes interiores de los
% parches infinitesimales se cancela y solo queda la del borde L = dS.
\begin{figure}[ht]
\centering
\begin{tikzpicture}[>=Stealth, x=1cm, y=1cm]
  \def\blob{plot[smooth cycle, tension=0.75] coordinates {(-4.3,0.1) (-3.1,1.25) (-1.2,1.5)
     (0.9,1.05) (2.7,1.45) (4.1,0.8) (4.5,-0.25) (3.3,-1.0) (1.0,-1.3) (-1.4,-1.1) (-3.3,-0.9)}}
  \path[name path=superficie] \blob;
  % superficie con degradado azul
  \shade[left color=azulclaro!35, right color=azulclaro!85] \blob;
  \begin{scope}
    \clip \blob;
    % parches con su circulación (elipses por la perspectiva)
    \foreach \j in {-3,-2,...,3}{\foreach \i in {-8,-7,...,8}{%
      \pgfmathsetmacro{\cx}{0.55*\i+0.12*\j}\pgfmathsetmacro{\cy}{0.33*\j-0.03*\i*\i/8}%
      \draw[black!70, line width=0.55pt,
            decoration={markings, mark=at position 0.2 with {\arrow{Stealth[length=2.6pt]}},
                                  mark=at position 0.7 with {\arrow{Stealth[length=2.6pt]}}},
            postaction={decorate}] (\cx,\cy) ellipse (0.27 and 0.145);}}
  \end{scope}
  % borde L con orientación
  \draw[rojoteo, line width=1.5pt,
        decoration={markings, mark=between positions 0.04 and 0.97 step 0.1 with {\arrow{Stealth[length=6pt]}}},
        postaction={decorate}] \blob;
  \draw[azulacero, line width=0.5pt, densely dashed] (0,0.1) -- ++(0,0.05);
  \draw[->, line width=1.6pt, azulacero] (0.15,0.25) -- ++(0,1.9) node[above, text=azulacero] {$\dS$};
  \node[font=\Large, text=azulprof] at (-3.4,0.1) {$S$};
  \node[font=\Large, text=rojoteo, anchor=west] at (4.6,-0.55) {$L=\partial S$};
\end{tikzpicture}
\caption{Idea del teorema de Stokes. La superficie $S$ se divide en parches
infinitesimales orientados, cada uno con su circulación (elipses). En las
paredes interiores cada lado se recorre dos veces en sentidos opuestos y la
circulación se cancela; solo sobrevive la circulación sobre el borde exterior
$L=\partial S$ (en rojo). La normal $\dS$ y el sentido de recorrido de $L$ están
ligados por la regla de la mano derecha.}
\label{fig:stokes}
\end{figure}


% -------------------------------------------------------------------------
\subsection{El teorema de la divergencia (de Gauss)}\label{sec:gauss}
% -------------------------------------------------------------------------
\begin{definicion}[Borde de un sólido y su orientación]\label{def:borde-solido}
El borde $\partial V$ de un sólido regular $V\subset\R^{3}$ es una o varias
superficies cerradas. La orientación positiva de $\partial V$ es la determinada
por el campo normal unitario \emph{exterior} (el que apunta fuera de $V$). Si $V$
tiene huecos o «burbujas», la normal exterior respecto a $V$ es la normal
\emph{interior} respecto a la burbuja.
\end{definicion}

\begin{teorema}[Teorema de la divergencia o de Gauss]\label{teo:gauss}
Sea $\vc{A}$ un campo vectorial de clase $C^{1}$ en un entorno de un sólido
acotado $V\subset\R^{3}$ cuyo borde $S=\partial V$ es una o varias superficies
cerradas orientadas con la normal exterior. Entonces
\begin{equation}\label{eq:gauss}
  \boxed{\;\iiint_{V}\nab\cdot\vc{A}\,\dV=\oiint_{\partial V}\vc{A}\cdot\dS\;}
\end{equation}
El flujo de $\vc{A}$ a través de la superficie cerrada que limita a $V$ es la
integral de su divergencia en el interior.
\end{teorema}
\begin{proof}[Idea de la demostración]
La fórmula es lineal en $\vc{A}$, así que basta probarla para campos paralelos
a los ejes: sea $\vc{A}=R\,\uz$. Para el cubo unidad $V=[0,1]^{3}$, solo las
tapas $z=1$ (normal $+\uz$) y $z=0$ (normal $-\uz$) contribuyen, pues en las
otras cuatro caras $\vc{A}\cdot\un=0$:
\begin{align*}
  \oiint_{\partial V}\vc{A}\cdot\dS&=\int_0^1\!\!\int_0^1\bigl[R(x,y,1)-R(x,y,0)\bigr]\dif x\,\dif y\\
  &=\int_0^1\!\!\int_0^1\!\!\int_0^1\pd{R}{z}\,\dif z\,\dif x\,\dif y=\iiint_V\nab\cdot(R\,\uz)\,\dV ,
\end{align*}
donde se ha usado el teorema fundamental del cálculo en $z$. Para un sólido
general se pega una malla de cubitos (o se cambia de variable): cada cara
interior aparece dos veces con normales opuestas, y los flujos interiores se
cancelan, de modo que solo queda el flujo a través del borde exterior
(\cref{fig:gauss-vol}).
\end{proof}

\begin{proposicion}[La divergencia como densidad de flujo]\label{prop:div-intrinseca}
Para un campo $\vc{F}$ de clase $C^{1}$ y un punto $\vc{p}$,
\begin{equation}\label{eq:div-intrinseca}
  (\nab\cdot\vc{F})(\vc{p})=\lim_{V\to\{\vc{p}\}}\frac{1}{v(V)}\oiint_{\partial V}\vc{F}\cdot\un\,\dsc ,
\end{equation}
siendo $V$ un pequeño sólido (por ejemplo una bola centrada en $\vc{p}$) de volumen
$v(V)$ que se contrae a $\vc{p}$.
\end{proposicion}
\begin{proof}
Por el \cref{teo:gauss}, $\frac{1}{v(V)}\oiint_{\partial V}\vc{F}\cdot\un\,\dsc=\frac{1}{v(V)}\iiint_V\nab\cdot\vc{F}\,\dV$,
que es el valor medio de la divergencia sobre $V$ y converge a $(\nab\cdot\vc{F})(\vc{p})$
al contraer $V$ (continuidad).
\end{proof}

La \cref{eq:div-intrinseca} es la definición general de la divergencia,
independiente del sistema de coordenadas, y explica su nombre y su
interpretación: es el flujo neto que sale por unidad de volumen. El teorema de
Gauss es entonces la expresión matemática de las cancelaciones de flujo en las
paredes interiores: la integral de la divergencia es la «suma» de todas las
fuentes y sumideros del interior, y solo sobrevive el flujo por el borde. Un
campo de divergencia nula (solenoidal, o incompresible) tiene flujo neto nulo a
través de toda superficie cerrada.

% Figura: teorema de la divergencia. Campo radial y elemento de volumen en
% coordenadas esféricas.
% Elemento de volumen esférico: caras como polígonos con aristas muestreadas.
% \pt{r}{theta}{phi} define las coordenadas cartesianas en \X, \Y, \Z
\newcommand{\sph}[3]{({(#1)*sin(#2)*cos(#3)},{(#1)*sin(#2)*sin(#3)},{(#1)*cos(#2)})}
% cara a r fijo: #1 = r, recorre theta (ta..tb) y phi (pa..pb)
\newcommand{\caraR}[1]{%
  \fill[ocreteo, opacity=0.38]
  ({#1*sin(\ta)*cos(\pa)},{#1*sin(\ta)*sin(\pa)},{#1*cos(\ta)})
  \foreach \k in {1,...,6}{ -- ({#1*sin(\ta+\k*(\tb-\ta)/6)*cos(\pa)},{#1*sin(\ta+\k*(\tb-\ta)/6)*sin(\pa)},{#1*cos(\ta+\k*(\tb-\ta)/6)}) }
  \foreach \k in {1,...,6}{ -- ({#1*sin(\tb)*cos(\pa+\k*(\pb-\pa)/6)},{#1*sin(\tb)*sin(\pa+\k*(\pb-\pa)/6)},{#1*cos(\tb)}) }
  \foreach \k in {1,...,6}{ -- ({#1*sin(\tb-\k*(\tb-\ta)/6)*cos(\pb)},{#1*sin(\tb-\k*(\tb-\ta)/6)*sin(\pb)},{#1*cos(\tb-\k*(\tb-\ta)/6)}) }
  \foreach \k in {1,...,6}{ -- ({#1*sin(\ta)*cos(\pb-\k*(\pb-\pa)/6)},{#1*sin(\ta)*sin(\pb-\k*(\pb-\pa)/6)},{#1*cos(\ta)}) }
  -- cycle;}
% cara a theta fijo: #1 = theta
\newcommand{\caraT}[1]{%
  \fill[ocreteo, opacity=0.38]
  ({\ra*sin(#1)*cos(\pa)},{\ra*sin(#1)*sin(\pa)},{\ra*cos(#1)})
  -- ({\rb*sin(#1)*cos(\pa)},{\rb*sin(#1)*sin(\pa)},{\rb*cos(#1)})
  \foreach \k in {1,...,6}{ -- ({\rb*sin(#1)*cos(\pa+\k*(\pb-\pa)/6)},{\rb*sin(#1)*sin(\pa+\k*(\pb-\pa)/6)},{\rb*cos(#1)}) }
  -- ({\ra*sin(#1)*cos(\pb)},{\ra*sin(#1)*sin(\pb)},{\ra*cos(#1)})
  \foreach \k in {1,...,6}{ -- ({\ra*sin(#1)*cos(\pb-\k*(\pb-\pa)/6)},{\ra*sin(#1)*sin(\pb-\k*(\pb-\pa)/6)},{\ra*cos(#1)}) }
  -- cycle;}
% cara a phi fijo: #1 = phi
\newcommand{\caraP}[1]{%
  \fill[ocreteo, opacity=0.38]
  ({\ra*sin(\ta)*cos(#1)},{\ra*sin(\ta)*sin(#1)},{\ra*cos(\ta)})
  -- ({\rb*sin(\ta)*cos(#1)},{\rb*sin(\ta)*sin(#1)},{\rb*cos(\ta)})
  \foreach \k in {1,...,6}{ -- ({\rb*sin(\ta+\k*(\tb-\ta)/6)*cos(#1)},{\rb*sin(\ta+\k*(\tb-\ta)/6)*sin(#1)},{\rb*cos(\ta+\k*(\tb-\ta)/6)}) }
  -- ({\ra*sin(\tb)*cos(#1)},{\ra*sin(\tb)*sin(#1)},{\ra*cos(\tb)})
  \foreach \k in {1,...,6}{ -- ({\ra*sin(\tb-\k*(\tb-\ta)/6)*cos(#1)},{\ra*sin(\tb-\k*(\tb-\ta)/6)*sin(#1)},{\ra*cos(\tb-\k*(\tb-\ta)/6)}) }
  -- cycle;}
% arista a (r,theta) fijos recorriendo phi, etc.
\newcommand{\aristaP}[2]{\draw[rojoteo, dashed, line width=0.5pt]
  ({#1*sin(#2)*cos(\pa)},{#1*sin(#2)*sin(\pa)},{#1*cos(#2)})
  \foreach \k in {1,...,6}{ -- ({#1*sin(#2)*cos(\pa+\k*(\pb-\pa)/6)},{#1*sin(#2)*sin(\pa+\k*(\pb-\pa)/6)},{#1*cos(#2)}) };}
\newcommand{\aristaT}[2]{\draw[rojoteo, dashed, line width=0.5pt]
  ({#1*sin(\ta)*cos(#2)},{#1*sin(\ta)*sin(#2)},{#1*cos(\ta)})
  \foreach \k in {1,...,6}{ -- ({#1*sin(\ta+\k*(\tb-\ta)/6)*cos(#2)},{#1*sin(\ta+\k*(\tb-\ta)/6)*sin(#2)},{#1*cos(\ta+\k*(\tb-\ta)/6)}) };}
\begin{figure}[ht]
\centering
\tdplotsetmaincoords{66}{128}
\begin{minipage}[t]{0.48\textwidth}\centering
\begin{tikzpicture}[tdplot_main_coords, scale=0.78, baseline=(current bounding box.north)]
  \draw[eje] (0,0,0) -- (3.3,0,0) node[below left] {\small $x$};
  \draw[eje] (0,0,0) -- (0,3.3,0) node[right] {\small $y$};
  \draw[eje] (0,0,0) -- (0,0,3.3) node[above] {\small $z$};
  % campo radial (malla gruesa, módulo ~ 1/r^0: flechas unitarias coloreadas por r)
  \foreach \z in {-1.8,-0.6,0.6,1.8}{\foreach \y in {1.8,0.6,-0.6,-1.8}{\foreach \x in {-1.8,-0.6,0.6,1.8}{%
    \pgfmathsetmacro{\fm}{sqrt(\x*\x+\y*\y+\z*\z)}%
    \pgfmathsetmacro{\ft}{\fm/3.12}\mapcolor{\ft}%
    \draw[-{Stealth[length=3.4pt,width=3.4pt]}, line width=0.7pt, mapc]
      (\x,\y,\z) -- ++({0.36*\x/\fm*1.6},{0.36*\y/\fm*1.6},{0.36*\z/\fm*1.6});}}}
  % elemento de volumen: r in [1.4,2.5], theta in [38,62], phi in [8,42]
  \def\ra{1.4}\def\rb{2.5}\def\ta{38}\def\tb{62}\def\pa{8}\def\pb{42}
  \caraR{\ra}\caraR{\rb}
  \caraT{\ta}\caraT{\tb}
  \caraP{\pa}\caraP{\pb}
  \aristaP{\ra}{\ta}\aristaP{\ra}{\tb}\aristaP{\rb}{\ta}\aristaP{\rb}{\tb}
  \aristaT{\ra}{\pa}\aristaT{\ra}{\pb}\aristaT{\rb}{\pa}\aristaT{\rb}{\pb}
  \foreach \tt in {\ta,\tb}{\foreach \pp in {\pa,\pb}{\draw[rojoteo, dashed, line width=0.5pt]
      \sph{\ra}{\tt}{\pp} -- \sph{\rb}{\tt}{\pp};}}
\end{tikzpicture}\\[2pt]{\small\scshape (a)}
\end{minipage}\hfill
\begin{minipage}[t]{0.48\textwidth}\centering
\begin{tikzpicture}[>=Stealth, x=1cm, y=1cm, baseline=(current bounding box.north)]
  \foreach \a in {0,18,...,342}{%
    \draw[-{Stealth[length=5pt,width=5pt]}, azulacero, line width=1pt]
      (\a:0.55) -- (\a:1.65);
    \draw[-{Stealth[length=5pt,width=5pt]}, azulacero, line width=1pt]
      (\a:1.95) -- (\a:2.55);
    \draw[-{Stealth[length=5pt,width=5pt]}, azulacero!70, line width=1pt]
      (\a:2.85) -- (\a:3.2);}
  \fill[azulprof] (0,0) circle (1.6pt);
  % elemento de área (corte del elemento de volumen): sector anular
  \draw[rojoteo, dashed, line width=0.9pt, fill=ocreteo, fill opacity=0.25]
     (60:1.7) arc (60:84:1.7) -- (84:2.9) arc (84:60:2.9) -- cycle;
  \draw[rojoteo, line width=0.4pt] (0,0) -- (60:1.7);
  \node[font=\footnotesize, text=rojoteo] at (72:3.35) {$r\,\dif\theta$};
  \node[font=\footnotesize, text=rojoteo, anchor=west] at (50:2.35) {$\dif r$};
\end{tikzpicture}\\[2pt]{\small\scshape (b)}
\end{minipage}\\[6pt]
\begin{minipage}[t]{0.95\textwidth}\centering
\begin{tikzpicture}[>=Stealth, x=1cm, y=1cm, baseline=(current bounding box.north)]
  % elemento de volumen (paralelepípedo en proyección oblicua)
  \def\ax{2.8}\def\ay{1.7}\def\ox{1.0}\def\oy{0.55}
  \draw[rojoteo, dashed, line width=0.7pt] (\ox,\oy) -- (\ox+\ax,\oy) (\ox,\oy) -- (\ox,\oy+\ay) (\ox,\oy) -- (0,0);
  \fill[ocreteo, opacity=0.28] (0,0) rectangle (\ax,\ay);
  \draw[rojoteo, line width=0.9pt] (0,0) rectangle (\ax,\ay);
  \fill[ocreteo, opacity=0.18] (0,\ay) -- (\ax,\ay) -- (\ax+\ox,\ay+\oy) -- (\ox,\ay+\oy) -- cycle;
  \fill[ocreteo, opacity=0.40] (\ax,0) -- (\ax,\ay) -- (\ax+\ox,\ay+\oy) -- (\ax+\ox,\oy) -- cycle;
  \draw[rojoteo, line width=0.9pt] (0,\ay) -- (\ox,\ay+\oy) -- (\ax+\ox,\ay+\oy) -- (\ax+\ox,\oy) -- (\ax,0)
        (\ax,\ay) -- (\ax+\ox,\ay+\oy);
  % terna de versores locales
  \coordinate (T) at (-3.7,1.0);
  \draw[->, line width=1.2pt, verdeteo] (T) -- ++(0.95,0.52) node[above right, font=\small] {$\ur$};
  \draw[->, line width=1.2pt, verdeteo] (T) -- ++(0,-1.15) node[below, font=\small] {$\uth$};
  \draw[->, line width=1.2pt, verdeteo] (T) -- ++(1.15,0) node[right, font=\small] {$\uph$};
  \fill[verdeteo] (T) circle (1.3pt);
  % cotas
  \node[font=\small, text=rojoteo, anchor=north] at (\ax/2,-0.08) {$r\sin\theta\,\dif\varphi$};
  \node[font=\small, text=rojoteo, anchor=east] at (-0.08,\ay/2) {$r\,\dif\theta$};
  \node[font=\small, text=rojoteo, anchor=south west] at (\ax+0.45,\ay+0.5) {$\dif r$};
  \node[font=\small, text=black!80, anchor=west] at (\ax+\ox+0.5,\ay/2+\oy/2) {$\dV=r^{2}\sin\theta\,\dif r\,\dif\theta\,\dif\varphi$};
\end{tikzpicture}\\[2pt]{\small\scshape (c)}
\end{minipage}
\caption{Teorema de la divergencia con simetría radial. (a)~Campo radial
$\vc{A}=A(r)\,\ur$ y un elemento de volumen $\dV$ limitado por
$r,r+\dif r$; $\theta,\theta+\dif\theta$; $\varphi,\varphi+\dif\varphi$.
(b)~Corte meridiano: el elemento es un sector anular de lados $\dif r$ y
$r\,\dif\theta$. (c)~El elemento, ampliado: es un «ladrillo» curvilíneo de aristas
$\dif r$, $r\,\dif\theta$ y $r\sin\theta\,\dif\varphi$ según la terna ortonormal
$(\ur,\uth,\uph)$. El flujo neto por sus seis caras es $\nab\cdot\vc{A}\,\dV$.}
\label{fig:gauss-vol}
\end{figure}


% -------------------------------------------------------------------------
\subsection{Resumen y comprobación en campos sencillos}\label{sec:resumen-teoremas}
% -------------------------------------------------------------------------
Los tres teoremas, en la notación que se empleará en el resto del curso, son:
\begin{align}
  \int_{\vc{a}}^{\vc{b}}\nab f\cdot\dl&=f(\vc{b})-f(\vc{a}) &&\text{(teorema del gradiente)},\label{eq:resumen-grad}\\
  \int_{V}\nab\cdot\vc{A}\,\dV&=\oint_{S}\vc{A}\cdot\dS &&\text{(teorema de la divergencia, Gauss)},\label{eq:resumen-gauss}\\
  \int_{S}(\nab\times\vc{A})\cdot\dS&=\oint_{L}\vc{A}\cdot\dl &&\text{(teorema del rotacional, Stokes)},\label{eq:resumen-stokes}
\end{align}
con $S=\partial V$ en la segunda y $L=\partial S$ en la tercera. Se reconoce en
todos ellos la estructura «integral de la derivada $=$ integral sobre el borde».

\begin{ejemplo}[Comprobación en los campos de referencia]\label{ej:comprobacion-campos}
Tomamos los campos de la \cref{fig:campos_div} y los dos objetos de la
\cref{fig:campos_int}: un cubo unidad $V=[1,2]\times[1{,}5,2{,}5]\times[0,1]$ y un
circuito rectangular $C=\partial S$ en el plano $z=0$ con
$S=[1,2]\times[0{,}5,1]$, recorrido en sentido horario (normal $-\uz$, de
área $0{,}5$).
\begin{itemize}
  \item $\vc{F}_2=x\,\ux$ ($\nab\cdot\vc{F}_2=1$, $\nab\times\vc{F}_2=\vc{0}$).
        Flujo por las tapas $x=2$ y $x=1$: $(2-1)\cdot1=1=\iiint_V1\,\dV$;
        el resto de caras no contribuye. Circulación: los lados horizontales
        se cancelan y los verticales son perpendiculares al campo, luego es
        nula, de acuerdo con $\nab\times\vc{F}_2=\vc{0}$.
  \item $\vc{F}_3=y\,\ux$ ($\nab\cdot\vc{F}_3=0$, $\nab\times\vc{F}_3=-\uz$).
        El flujo por el cubo es nulo (lo que entra por una cara sale por la
        opuesta), y la circulación horaria es
        $\oint\vc{F}_3\cdot\dl=(1-0{,}5)\cdot1=0{,}5=\iint_S(-\uz)\cdot(-\uz)\,\dsc$
        (el lado superior, a $y=1$, está a favor del campo, el inferior, a
        $y=0{,}5$, en contra y más débil).
\end{itemize}
\end{ejemplo}

% Figura: los campos de referencia con un volumen V (cubo) y un circuito C.
\newcommand{\cuboV}[2]{%  esquina inferior izquierda (#1,#2): cubo de lado 1
  \fill[grisclaro] (#1,#2) rectangle (#1+1,#2+1);
  \fill[grispizarra!40] (#1,#2+1) -- (#1+0.35,#2+1.3) -- (#1+1.35,#2+1.3) -- (#1+1,#2+1) -- cycle;
  \fill[grispizarra!55] (#1+1,#2) -- (#1+1.35,#2+0.3) -- (#1+1.35,#2+1.3) -- (#1+1,#2+1) -- cycle;
  \draw[rojoteo, line width=0.8pt] (#1,#2) rectangle (#1+1,#2+1);
  \draw[rojoteo, line width=0.8pt] (#1,#2+1) -- (#1+0.35,#2+1.3) -- (#1+1.35,#2+1.3) -- (#1+1.35,#2+0.3) -- (#1+1,#2);
  \draw[rojoteo, line width=0.8pt] (#1+1,#2+1) -- (#1+1.35,#2+1.3);
  \node[font=\footnotesize, text=rojoteo] at (#1+0.5,#2+0.5) {$V$};}
\newcommand{\circuitoC}[4]{%  rectángulo [#1,#2]x[#3,#4] recorrido en sentido horario
  \fill[white, opacity=0.55] (#1,#3) rectangle (#2,#4);\draw[rojoteo, line width=1pt] (#1,#3) rectangle (#2,#4);
  \draw[-{Stealth[length=6pt]}, rojoteo, line width=1pt] ({(#1+#2)/2-0.15},#4) -- ({(#1+#2)/2+0.05},#4);
  \draw[-{Stealth[length=6pt]}, rojoteo, line width=1pt] (#2,{(#3+#4)/2+0.12}) -- (#2,{(#3+#4)/2-0.05});
  \draw[-{Stealth[length=6pt]}, rojoteo, line width=1pt] ({(#1+#2)/2+0.15},#3) -- ({(#1+#2)/2-0.05},#3);
  \draw[-{Stealth[length=6pt]}, rojoteo, line width=1pt] (#1,{(#3+#4)/2-0.12}) -- (#1,{(#3+#4)/2+0.05});
  \node[font=\footnotesize, text=rojoteo, anchor=west] at (#2+0.05,{(#3+#4)/2}) {$C$};}
\begin{figure}[ht]
\centering
\begin{minipage}[t]{0.325\textwidth}\centering
\begin{tikzpicture}[x=0.68cm, y=0.68cm, baseline=(current bounding box.north)]
  \marcoD{-2.7}{2.7}{-2.7}{2.7}{-2,-1,0,1,2}{-2,-1,0,1,2}
  \ejesD{-2.7}{2.7}{-2.7}{2.7}
  \foreach \x in {-2.4,-1.8,...,2.4}{\foreach \y in {-2.4,-1.8,...,2.4}{%
    \draw[-{Stealth[length=4.5pt,width=4.5pt]}, line width=0.7pt, azulprof]
      (\x-0.17,\y-0.17) -- ++(0.34,0.34);}}
\end{tikzpicture}\\[3pt]{\small $\vc{F}_1=\ux+\uy$}
\end{minipage}\hfill
\begin{minipage}[t]{0.325\textwidth}\centering
\begin{tikzpicture}[x=1.12cm, y=1.12cm, baseline=(current bounding box.north)]
  \marcoD{0}{3.2}{0}{3.2}{0,1,2,3}{0,1,2,3}
  \ejesD{0}{3.2}{0}{3.2}
  \clipD{0}{3.2}{0}{3.2}{\flechas{\x}{0}{0.2,0.575,...,3.1}{0.2,0.575,...,3.1}{0.30}{3.1}}
  \cuboV{1}{1.5}
  \circuitoC{1}{2}{0.5}{1}
\end{tikzpicture}\\[3pt]{\small $\vc{F}_2=x\,\ux$}
\end{minipage}\hfill
\begin{minipage}[t]{0.325\textwidth}\centering
\begin{tikzpicture}[x=1.12cm, y=1.12cm, baseline=(current bounding box.north)]
  \marcoD{0}{3.2}{0}{3.2}{0,1,2,3}{0,1,2,3}
  \ejesD{0}{3.2}{0}{3.2}
  \clipD{0}{3.2}{0}{3.2}{\flechas{\y}{0}{0.2,0.575,...,3.1}{0.2,0.575,...,3.1}{0.30}{3.1}}
  \cuboV{1}{1.5}
  \circuitoC{1}{2}{0.5}{1}
\end{tikzpicture}\\[3pt]{\small $\vc{F}_3=y\,\ux$}
\end{minipage}
\caption{Campos de referencia con un volumen $V$ (cubo) y un circuito $C$
(rectángulo recorrido en sentido horario). Para $\vc{F}_2$ el flujo neto por la
superficie del cubo es no nulo ($=1$) y la circulación por $C$ es nula; para
$\vc{F}_3$ ocurre lo contrario: el flujo es nulo y la circulación vale $0{,}5$.}
\label{fig:campos_int}
\end{figure}


\begin{ejemplo}[Rotacional uniforme: la circulación solo depende del área]\label{ej:rot-uniforme}
Los campos $\vc{F}=y\,\ux-x\,\uy$ y $\vc{F}=y\,\ux$ tienen rotacional
\emph{uniforme} (el mismo en todos los puntos): $-2\,\uz$ y $-\uz$,
respectivamente. Por el teorema de Stokes, la circulación horaria (normal $-\uz$)
alrededor de un circuito de área $a$ vale $\oint\vc{F}\cdot\dl=\iint(\nab\times\vc{F})\cdot(-\uz)\,\dsc=2a$
(resp. $a$), \emph{sin importar dónde esté colocado el circuito}. En la
\cref{fig:campos_circ}, los dos rectángulos de área $a=1{,}4$ dan $2{,}8$ (resp.
$1{,}4$) aunque el campo valga cosas muy distintas sobre cada uno. Se comprueba
directamente en el segundo caso: el campo es horizontal, los lados verticales
no aportan, y los horizontales dan $(y_{\mathrm{sup}}-y_{\mathrm{inf}})\cdot1{,}4=1\cdot1{,}4=1{,}4$,
con independencia de la altura $y$ a la que esté el rectángulo.
\end{ejemplo}

% Figura: campos de rotacional uniforme y dos circuitos iguales en posiciones distintas.
\begin{figure}[ht]
\centering
\begin{minipage}[t]{0.47\textwidth}\centering
\begin{tikzpicture}[x=0.74cm, y=0.74cm, baseline=(current bounding box.north)]
  \marcoD{-2.7}{2.7}{-2.7}{2.7}{-2,-1,0,1,2}{-2,-1,0,1,2}
  \ejesD{-2.7}{2.7}{-2.7}{2.7}
  \clipD{-2.7}{2.7}{-2.7}{2.7}{\flechas{\y}{-\x}{-2.4,-1.8,...,2.4}{-2.4,-1.8,...,2.4}{0.34}{3.4}}
  \circuitoC{-0.7}{0.7}{-0.45}{0.55}
  \circuitoC{-0.7}{0.7}{-2.0}{-1.0}
\end{tikzpicture}\\[3pt]{\small $\vc{F}=y\,\ux-x\,\uy$,\quad $\nab\times\vc{F}=-2\,\uz$}
\end{minipage}\hfill
\begin{minipage}[t]{0.47\textwidth}\centering
\begin{tikzpicture}[x=0.74cm, y=0.74cm, baseline=(current bounding box.north)]
  \marcoD{-2.7}{2.7}{-2.7}{2.7}{-2,-1,0,1,2}{-2,-1,0,1,2}
  \ejesD{-2.7}{2.7}{-2.7}{2.7}
  \clipD{-2.7}{2.7}{-2.7}{2.7}{\flechas{\y}{0}{-2.4,-1.8,...,2.4}{-2.4,-1.8,...,2.4}{0.34}{2.5}}
  \circuitoC{-0.7}{0.7}{-0.45}{0.55}
  \circuitoC{-0.7}{0.7}{-2.0}{-1.0}
\end{tikzpicture}\\[3pt]{\small $\vc{F}=y\,\ux$,\quad $\nab\times\vc{F}=-\uz$}
\end{minipage}
\caption{Campos de rotacional uniforme y dos circuitos rectangulares iguales
(área $1{,}4$, recorridos en sentido horario) en posiciones distintas. Aunque
el campo es muy diferente sobre uno y otro circuito, la circulación es la misma
en ambos: $2{,}8$ para $\vc{F}=y\,\ux-x\,\uy$ y $1{,}4$ para $\vc{F}=y\,\ux$.
Solo depende del rotacional y del área encerrada.}
\label{fig:campos_circ}
\end{figure}


